\documentclass[12pt]{article}
\usepackage[utf8]{inputenc}
\usepackage[T1]{fontenc}
\usepackage{hyperref}
\title{Kwegu Language Resources for K-12 Curriculum Development}
\author{Compiled by Claude Sonnet 5.5 (claude-sonnet-5-5) for Andr\'as Kornai}
\date{}
\begin{document}
\maketitle

\section{Introduction}
Kwegu (ISO 639-3 xwg; also Koegu, Bacha, Kwegi, Menja, Nidi) is a Surmic language spoken on the west bank of the Omo River in southwestern Ethiopia. Speaker counts vary: about 300 in a 1987--1990 count, 3,900 (Ethnologue 2016), and 5,100 ethnic Kwegu in the 2007 census. Hieda describes its lexicon and syntax as unusual within Surmic because of contact with the Omotic language Kara. Documentation exists in academic sources, but no school materials were found.

\section{Primary Educational Resources}
No primers, textbooks or school programmes for Kwegu were found in the searches.

\section{Linguistic and Reference Materials}
\begin{itemize}
\item Hieda's Koegu vocabulary study, Kyoto University repository: \url{https://repository.kulib.kyoto-u.ac.jp/dspace/bitstream/2433/68358/1/ASM_S_14_1.pdf}
\item Osaka University repository item (Hieda; the PDF was listed but its content was not checked): \url{https://ir.library.osaka-u.ac.jp/repo/ouka/all/71074/saak_003_131.pdf}
\item Wikipedia, Kwegu language: \url{https://en.wikipedia.org/wiki/Kwegu_language}
\item Wikipedia, Kwegu people: \url{https://en.wikipedia.org/wiki/Kwegu_people}
\item Endangered Languages Project record: \url{https://endangeredlanguages.com/lang/2677?hl=en}
\end{itemize}
The search summary mentions an ELAR archive of Koegu materials, but no link was seen, so it is not listed.

\section{Religious and Community Resources}
Joshua Project profile: \url{https://joshuaproject.net/people_groups/18363/ET}. No scripture or community organisation was identified in the results.

\section{Digital and Audio Resources}
No audio, video or software for Kwegu was found in the searches.

\section{Curriculum Development Context}
Kwegu is endangered. The people are small-scale cultivators and hunters living among larger neighbouring groups, and no school use is known. Linguistic documentation is mainly academic (Hieda's work). A curriculum would first need community consultation, an orthography, and collaboration with documentation holders. The ELAR archive could be explored for audio and texts.
\end{document}